\documentclass[10pt,a4paper]{amsart}
\usepackage[margin=19mm]{geometry}
\usepackage{amsmath,amssymb,mathtools}
\usepackage[scr=rsfs,cal=euler]{mathalfa}
\usepackage{xfrac}
\usepackage[unicode,hidelinks,bookmarks=false]{hyperref}
\newcommand{\ie}{i.e.\ }
\newcommand{\cf}{cf.\ }
\newcommand{\resp}{respectively\ }
\newcommand{\Subsection}{\subsection}
\providecommand{\nobreakdash}{\nobreak\hbox{-}}
\setlength{\emergencystretch}{2em}
\multlinegap0pt
\usepackage{amssymb}
\usepackage{enumerate}
\usepackage{algorithm}\usepackage{algpseudocode}
\newtheorem{proposition}{Proposition}
\newcommand{\F}{\mathbb F}
\newcommand{\G}{\Gamma}
\title{Graphs from quadratic forms and vector spaces over finite fields}
\author{Jean Godard, Lucas Reis}
\date{Source: arXiv:2605.21866v1 (CC BY 4.0)}
\begin{document}
\maketitle

\noindent\emph{Source and licence.} Graphs from quadratic forms and vector spaces over finite fields. Version arXiv:2605.21866v1 (CC BY 4.0), distributed by arXiv under the Creative Commons Attribution 4.0 International licence, http://creativecommons.org/licenses/by/4.0/, as recorded in the arXiv licence metadata for this exact version and shown on the abstract page https://arxiv.org/abs/2605.21866v1. Full licence notice: \texttt{licenses/arxiv-2605.21866-CC-BY-4.0.txt}.\par\smallskip

\noindent\textit{Editorial scope.} Counted unit: the article's proposition prop:Q+ (source0.tex lines 259--261). The article's complete original proof of it is delivered with it (source0.tex lines 262--270, one \texttt{proof} environment of 9 lines). No mathematics is written, repaired or re-derived: the statement and the proof are the article's own lines, in the article's own notation, and no retained line is substituted or re-flowed.  No further source-local context is needed: every cross-reference of every retained line resolves inside the two delivered spans.  Nothing was omitted from any retained span, so the package carries no omission record.  No retained line contains a citation, so the wrapper carries no bibliography and no bibliography entry.  No retained line contains a figure environment, an image-inclusion command or drawing code, so no image is delivered and the package declares no asset dependency.  Editorial formatting only, in addition: an AMS \texttt{amsart} preamble that restates the article's own declarations and macros used by the retained lines, namely the environments \texttt{proposition} on separate counters, together with the standard packages those lines need.  The retained environments print in this document as Proposition 1 in order of appearance, whereas the article prints the same items with the numbering of its own full document.  The complete native source of the article as distributed in the arXiv e-print for this exact version is bundled unchanged as \texttt{sources/201093/source0.tex}.
\begin{proposition}\label{prop:Q+}
    Let $V\subsetneq \F_{q^n}$ be a proper $\F_q$-vector subspace and let $S_V$ be the set of squares in $V$. If $S_V=\{0\}$, then $\omega(Q_+, V)=2$ and $\G(Q_+, V)$ is triangle-free. Otherwise, the set $C_V:=\{u\in \F_{q^n}: u^2\in S_V\}$ induces a maximal clique in $\Gamma(Q_+, V)$ of size at least $3$, and every other maximal clique of this graph has size $2$. In particular, $\Gamma(Q_+, V)$ is disconnected.
\end{proposition}

\begin{proof}
For the first statement, suppose that $\omega(Q_+, V)\ge 3$ and let $x_1, x_2, x_3\in \F_{q^n}$ be elements inducing a triangle in $\G(Q_+, V)$ with $x_1\ne 0$. Hence $x_1^2+x_2^2, x_1^2+x_3^2, x_2^2+x_3^2\in V$. Since $V$ is a vector subspace, it follows that
$$2x_1^2=(x_1^2+x_2^2)+(x_1^2+x_3^2)-(x_2^2+x_3^2)\in V.$$
Since $q$ is odd, we obtain $x_1^2\in V$ and then $S_V\ne \{0\}$.
For the second statement, suppose $S_V\ne \{0\}$ and let $u^2\in S_V$ with $u\ne 0$. It is direct to verify that $C_V$ induces a clique containing the vertices $u, -u$ and $0$. On the other hand, from the proof of the first statement, the vertices of every clique of size at least $3$ must lie in $C_V$. This completes the proof of the second statement. 

For the last statement, observe that if $(x, v)$ is an edge of $\G(Q_+, V)$
 with $v\in C_V$, then $x^2+v^2\in V$. Since $v^2\in V$, we obtain $x^2\in V$ and so $x\in C_V$. Hence $C_V$ induces a maximal connected component of $\G(Q_+, V)$. It is clear that $C_V$ has size at most $2\# V$. Therefore, if $\G(Q_+, V)$ is connected, then $2\# V\ge q^n$ and, since $q$ is odd, the latter is equivalent to $\# V=q^n$, a contradiction with $V\ne \F_{q^n}$.
\end{proof}

\end{document}
